\scanpage{24}
\mspage{\textcircled{6}}

\begin{center}
\uline{ACADEMIE D'AIX-MARSEILLE}\\
INSPECTION ACADEMIQUE\\
\uline{DE VAUCLUSE}
\end{center}

Quelle est la partie de $X_2\times X_2$ qui provient de $X'_3$ ?

\[
\text{a)}\quad(p,p\rho)(e_3)=(e_2,e_2)
\]
\[
\text{b)}\quad(p,p\rho)\underbrace{\rho^i d(\unsure{\xi|\zeta})}_{\substack{i\in\mathbb Z/3\mathbb Z\\\unsure{\xi|\zeta}\in X_2^*=X_2\setminus\{e_2\}}}
=\begin{cases}
(\unsure{\xi|\zeta},e_2)&\text{si }i=0\\
(e_2,\sigma\unsure{\xi|\zeta})&\text{si }i=1\\
(\sigma\unsure{\xi|\zeta},\unsure{\xi|\zeta})&\text{si }i=2
\end{cases}
\]

Donc on trouve les éléments \deleted{de la forme} dans
\[
\uline{\underbrace{(X_2\times\{e_2\})}_{\Sigma_1}\cup\underbrace{(\{e_2\}\times X_2)}_{\Sigma_2}\cup\deleted{\unsure{\Delta}}\Delta'}
\]
où $\Delta'$ est l'image de
\diagram{\begin{tikzcd}[ampersand replacement=\&]
X_2\arrow[r,"{\delta'=(\sigma,\mathrm{id}_{X_2})}"']\&X_2\times X_2
\end{tikzcd}}

\uline{NB} Cette partie $\Sigma$ de $X_2\times X_2$ est la réunion de $\Sigma_1,\Sigma_2,\Delta'$, et l'intersection de deux \unsure{quelconques} des trois est égale à $\{(e_2,e_2)\}$.

Quand a-t-on $\Sigma=X_2\times X_2$ ?

\deleted{\unsure{Si}} On a
\[
X_2\times X_2\setminus\{\Sigma_1\cup\Sigma_2)=X_2^*\times X_2^*,
\]
donc il est \unsure{nécs.} que $X_2^*\longrightarrow X_2^*\times X_2^*$ soit surj., i.e. que pour $\unsure{\zeta|\xi},\eta\in X_2^*$, on ait tjs $\eta=\sigma\unsure{\zeta|\xi}$. Cela signifie que l'on est dans l'un des 2 cas suivants

a) $X_2^*=\varnothing$ i.e. $X_2=\{e_2\}$

b) $X_2^*$ de cardinal 1 (donc $\sigma$ \deleted{\unsure{y}} y opère \unsure{trivialt})

\margin{i.e. $X_2$ de cardinal 2, (et $\sigma$ y opère \unsure{trivialt})}

[Si $X_2^*$ est de cardinal $\geq2$, alors s'il y a une seule orbite sous $\sigma$, celle-ci est de cardinal 2, et \deleted{$\unsure{\xi}$} $\forall\unsure{\zeta|\xi}\in X_2^*$, $(\unsure{\zeta|\xi},\unsure{\zeta|\xi})\notin\operatorname{Im}\delta'$ i.e. $\unsure{\zeta|\xi}\ne\sigma\unsure{\zeta|\xi}$). S'il y a \deleted{\unsure{donc}} au moins deux orbites, on prend $\unsure{\zeta|\xi},\eta$ dans deux orbites différentes, et on a bien $\eta\ne\sigma\unsure{\zeta|\xi}$.

\margin{23}
